% ============================================================
%  Chapter 4 -- Continuous Functions, 4.5-4.12
% ============================================================
\rsection{Continuous Functions}

\begin{signpost}[Three definitions converging]
Continuity arrives here in its $\varepsilon$--$\delta$ form (4.5), is linked
to limits (4.6) and hence to sequences (through 4.2), and is then recast as
the inverse-image condition (4.8). All three are the same property; each has
its use. $\varepsilon$--$\delta$ verifies concrete examples (4.11);
sequences refute; 4.8 proves theorems --- it is the only form in which the
compactness and connectedness arguments of the next two sections can even be
stated. The section's quiet workhorses are 4.7 (compositions) and the
assembly line of 4.9--4.11, which manufactures the continuity of every
rational function once and for all.
\end{signpost}

\ritem{4.5}{Definition}
Suppose $X$ and $Y$ are metric spaces, $E \subset X$, $p \in E$, and $f$ maps
$E$ into $Y$. Then $f$ is said to be \emph{continuous at $p$} if for every
$\varepsilon > 0$ there exists a $\delta > 0$ such that
\[ d_Y(f(x),f(p)) < \varepsilon \]
for all points $x \in E$ for which $d_X(x,p) < \delta$.

If $f$ is continuous at every point of $E$, then $f$ is said to be
\emph{continuous on $E$}.

It should be noted that $f$ has to be defined at the point $p$ in order to be
continuous at $p$. (Compare this with the remark following Definition
4.1.)\kw{The two definitions differ in exactly two characters: 4.5 requires
$p \in E$ where 4.1 required $p$ to be a limit point, and 4.5 drops the
$0 <$ from the distance condition --- the point $p$ itself is now tested, and
the target is $f(p)$, not a free $q$. Keep the ledger straight: limits
ignore $p$; continuity is \emph{about} $p$.}

If $p$ is an isolated point of $E$, then our definition implies that every
function $f$ which has $E$ as its domain of definition is continuous at $p$.
For, no matter which $\varepsilon > 0$ we choose, we can pick $\delta > 0$ so
that the only point $x \in E$ for which $d_X(x,p) < \delta$ is $x = p$; then
\[ d_Y(f(x),f(p)) = 0 < \varepsilon. \]\kn{Not a pathology but a consistency
check: at an isolated point there is no approach to make, so nothing to fail.
It is also why 4.6 needs its extra hypothesis --- at an isolated point the
limit in 4.6 is undefined (4.1 demands a limit point) while continuity holds
vacuously. Every sequence-in-$\Z$, for instance, is a continuous function.}

\ritem{4.6}{Theorem}
\emph{In the situation given in Definition 4.5, assume also that $p$ is a
limit point of $E$. Then $f$ is continuous at $p$ if and only if
$\lim_{x\to p} f(x) = f(p)$.}

\begin{proof}
This is clear if we compare Definitions 4.1 and 4.5.
\end{proof}

We now turn to compositions of functions. A brief statement of the following
theorem is that a continuous function of a continuous function is continuous.

\ritem{4.7}{Theorem}
\emph{Suppose $X$, $Y$, $Z$ are metric spaces, $E \subset X$, $f$ maps $E$
into $Y$, $g$ maps the range of $f$, $f(E)$, into $Z$, and $h$ is the mapping
of $E$ into $Z$ defined by}
\[ h(x) = g(f(x)) \qquad (x \in E). \]
\emph{If $f$ is continuous at a point $p \in E$ and if $g$ is continuous at
the point $f(p)$, then $h$ is continuous at $p$.}

This function $h$ is called the \emph{composition} or the \emph{composite} of
$f$ and $g$. The notation
\[ h = g \circ f \]
is frequently used in this context.

\begin{proof}
Let $\varepsilon > 0$ be given. Since $g$ is continuous at $f(p)$, there
exists $\eta > 0$ such that
\[ d_Z(g(y),g(f(p))) < \varepsilon \quad\text{if } d_Y(y,f(p)) < \eta
   \text{ and } y \in f(E). \]
Since $f$ is continuous at $p$, there exists $\delta > 0$ such that
\[ d_Y(f(x),f(p)) < \eta \quad\text{if } d_X(x,p) < \delta \text{ and }
   x \in E. \]
It follows that
\[ d_Z(h(x),h(p)) = d_Z(g(f(x)),g(f(p))) < \varepsilon \]
if $d_X(x,p) < \delta$ and $x \in E$. Thus $h$ is continuous at
$p$.\kn{A relay race, run backwards from the goal: the final tolerance
$\varepsilon$ is handed to $g$, which quotes its price $\eta$; that price
becomes the tolerance handed to $f$, which quotes $\delta$. Chains of
quantifiers always compose in this order --- last function first --- and the
middle quantity $\eta$ is bought from one contract and spent on the other.
The same relay reappears whenever composed processes must be estimated
(e.g.\ 7.12).}
\end{proof}

\ritem{4.8}{Theorem}
\emph{A mapping $f$ of a metric space $X$ into a metric space $Y$ is
continuous on $X$ if and only if $f^{-1}(V)$ is open in $X$ for every open
set $V$ in $Y$.}

(Inverse images are defined in Definition 2.2.) This is a very useful
characterization of continuity.

\begin{proof}
Suppose $f$ is continuous on $X$ and $V$ is an open set in $Y$. We have to
show that every point of $f^{-1}(V)$ is an interior point of $f^{-1}(V)$. So,
suppose $p \in X$ and $f(p) \in V$. Since $V$ is open, there exists
$\varepsilon > 0$ such that $y \in V$ if $d_Y(f(p),y) < \varepsilon$; and
since $f$ is continuous at $p$, there exists $\delta > 0$ such that
$d_Y(f(x),f(p)) < \varepsilon$ if $d_X(x,p) < \delta$. Thus
$x \in f^{-1}(V)$ as soon as $d_X(x,p) < \delta$.

Conversely, suppose $f^{-1}(V)$ is open in $X$ for every open set $V$ in $Y$.
Fix $p \in X$ and $\varepsilon > 0$, let $V$ be the set of all $y \in Y$ such
that $d_Y(y,f(p)) < \varepsilon$. Then $V$ is open; hence $f^{-1}(V)$ is
open; hence there exists $\delta > 0$ such that $x \in f^{-1}(V)$ as soon as
$d_X(p,x) < \delta$. But if $x \in f^{-1}(V)$, then $f(x) \in V$, so that
$d_Y(f(x),f(p)) < \varepsilon$.

This completes the proof.
\end{proof}

\begin{proofwalk}[How the proof flows]
  \item \textbf{See the two vocabularies.} Continuity speaks in
        $\varepsilon$--$\delta$; openness speaks in interior points, i.e.\ in
        balls. The theorem claims they say the same thing, so the proof must
        be a two-way translation, and each direction is found by matching
        ball to ball.
  \item \textbf{Forward.} To show $f^{-1}(V)$ open, stand at $p \in
        f^{-1}(V)$. Openness of $V$ supplies an $\varepsilon$-ball inside $V$
        around $f(p)$; continuity at $p$ converts it into a $\delta$-ball
        around $p$ that $f$ sends into it. That $\delta$-ball certifies $p$
        interior. Note the order: $V$ speaks first, continuity answers.
  \item \textbf{Backward.} To verify continuity at $p$ for a given
        $\varepsilon$, \emph{manufacture} the open set to pull back: the
        $\varepsilon$-ball $V$ about $f(p)$ --- open by 2.19. The hypothesis
        opens $f^{-1}(V)$, and any $\delta$-ball inside it around $p$ is
        exactly the $\delta$ that continuity demanded.
  \item \textbf{Audit.} Each direction uses its hypothesis once, on the
        ball handed over by the other side; 2.19 (balls are open) is the only
        outside fact. The theorem is a change of coordinates, not a
        computation --- and note it is stated for $f$ on all of $X$, using
        Remark 4.12's licence to forget ambient sets.
\end{proofwalk}

\begin{center}
\begin{tikzpicture}[x=1mm, y=1mm]
  % X side
  \draw[thin] plot[smooth cycle, tension=0.8]
    coordinates {(0,2) (22,-2) (34,6) (28,20) (6,20)};
  \node[lbl, anchor=south west] at (-2,20) {$X$};
  \fill[black!10] plot[smooth cycle, tension=0.8]
    coordinates {(8,5) (20,3) (24,11) (13,15)};
  \draw[guide] plot[smooth cycle, tension=0.8]
    coordinates {(8,5) (20,3) (24,11) (13,15)};
  \node[lbl] at (16,1) {$f^{-1}(V)$};
  \node[ptfill, minimum size=2.4pt] at (15,9) {};
  \draw[thin] (15,9) circle (3mm);
  \node[mlbl, anchor=south] at (15,12.5) {$\delta$};
  \node[mlbl, anchor=north east] at (14.6,8.6) {$p$};
  % arrow
  \draw[vec] (38,9) -- (50,9);
  \node[mlbl, anchor=south] at (44,10.5) {$f$};
  % Y side
  \draw[thin] plot[smooth cycle, tension=0.8]
    coordinates {(54,2) (76,-2) (88,6) (82,20) (60,20)};
  \node[lbl, anchor=south east] at (90,20) {$Y$};
  \fill[black!10] plot[smooth cycle, tension=0.8]
    coordinates {(62,5) (74,3) (78,11) (67,15)};
  \draw[guide] plot[smooth cycle, tension=0.8]
    coordinates {(62,5) (74,3) (78,11) (67,15)};
  \node[lbl] at (70,1) {$V$ open};
  \node[ptfill, minimum size=2.4pt] at (69,9) {};
  \draw[thin] (69,9) circle (3mm);
  \node[mlbl, anchor=south] at (69,12.5) {$\varepsilon$};
  \node[mlbl, anchor=north east] at (68.6,8.6) {$f(p)$};
\end{tikzpicture}
\end{center}
\figcap{Theorem 4.8: whatever ball fits inside $V$ at $f(p)$, continuity
supplies a ball at $p$ mapping into it --- so $f^{-1}(V)$ is all interior.
The dashed regions correspond under $f$; the balls are the currency
exchanged.}

\ritem{}{Corollary}
\emph{A mapping $f$ of a metric space $X$ into a metric space $Y$ is
continuous if and only if $f^{-1}(C)$ is closed in $X$ for every closed set
$C$ in $Y$.}

This follows from the theorem, since a set is closed if and only if its
complement is open, and since $f^{-1}(E^c) = [f^{-1}(E)]^c$ for every
$E \subset Y$.\kn{The complement identity is exactly the third displayed law
in the box at 2.2 --- inverse images respect complements, images do not. Here
is the promised payoff: no image-based version of 4.8 exists, and the
asymmetry noted in Chapter 2 becomes the reason this whole chapter runs on
$f^{-1}$.}

We now turn to complex-valued and vector-valued functions, and to functions
defined on subsets of $\R^k$.

\ritem{4.9}{Theorem}
\emph{Let $f$ and $g$ be complex continuous functions on a metric space $X$.
Then $f+g$, $fg$, and $f/g$ are continuous on $X$.}

In the last case, we must of course assume that $g(x) \ne 0$, for all
$x \in X$.

\begin{proof}
At isolated points of $X$ there is nothing to prove. At limit points, the
statement follows from Theorems 4.4 and 4.6.
\end{proof}

\ritem{4.10}{Theorem}
\begin{enumerate}[label=(\alph*), leftmargin=2.2em, itemsep=3pt, topsep=3pt]
  \item \emph{Let $f_1,\dots,f_k$ be real functions on a metric space $X$,
        and let $\mathbf{f}$ be the mapping of $X$ into $\R^k$ defined by}
        \begin{equation}
          \mathbf{f}(x) = (f_1(x),\dots,f_k(x)) \qquad (x \in X); \tag{4.7}
        \end{equation}
        \emph{then $\mathbf{f}$ is continuous if and only if each of the
        functions $f_1,\dots,f_k$ is continuous.}
  \item \emph{If $\mathbf{f}$ and $\mathbf{g}$ are continuous mappings of $X$
        into $\R^k$, then $\mathbf{f}+\mathbf{g}$ and
        $\mathbf{f}\cdot\mathbf{g}$ are continuous on $X$.}
\end{enumerate}
The functions $f_1,\dots,f_k$ are called the \emph{components} of
$\mathbf{f}$. Note that $\mathbf{f}+\mathbf{g}$ is a mapping into $\R^k$,
whereas $\mathbf{f}\cdot\mathbf{g}$ is a real function on $X$.

\begin{proof}
Part (a) follows from the inequalities
\[ |f_j(x)-f_j(y)| \le |\mathbf{f}(x)-\mathbf{f}(y)| =
   \left\{\sum_{i=1}^{k} |f_i(x)-f_i(y)|^2\right\}^{1/2}, \]
for $j = 1,\dots,k$. Part (b) follows from (a) and Theorem 4.9.\kn{The same
two-sided squeeze as 3.4(a): each component error is at most the vector
error, and the vector error is at most $\sqrt k$ times the worst component.
Coordinatewise reduction, third appearance --- it will carry Chapter 9's
entire differential calculus.}
\end{proof}

\ritem{4.11}{Examples}
If $x_1,\dots,x_k$ are the coordinates of the point $\mathbf{x} \in \R^k$,
the functions $\phi_i$ defined by
\begin{equation}
  \phi_i(\mathbf{x}) = x_i \qquad (\mathbf{x} \in \R^k) \tag{4.8}
\end{equation}
are continuous on $\R^k$, since the inequality
\[ |\phi_i(\mathbf{x})-\phi_i(\mathbf{y})| \le |\mathbf{x}-\mathbf{y}| \]
shows that we may take $\delta = \varepsilon$ in Definition 4.5.

The functions $\phi_i$ are sometimes called the \emph{coordinate functions}.

Repeated application of Theorem 4.9 then shows that every \emph{monomial}
\begin{equation}
  x_1^{n_1}x_2^{n_2}\cdots x_k^{n_k} \tag{4.9}
\end{equation}
where $n_1,\dots,n_k$ are nonnegative integers, is continuous on $\R^k$. The
same is true of constant multiples of (4.9), since constants are evidently
continuous. It follows that every \emph{polynomial} $P$, given by
\begin{equation}
  P(\mathbf{x}) = \sum c_{n_1\dots n_k}x_1^{n_1}\cdots x_k^{n_k}
  \qquad (\mathbf{x} \in \R^k), \tag{4.10}
\end{equation}
is continuous on $\R^k$. Here the coefficients $c_{n_1\dots n_k}$ are complex
numbers, $n_1,\dots,n_k$ are nonnegative integers, and the sum in (4.10) has
finitely many terms.

Furthermore, every \emph{rational function} in $x_1,\dots,x_k$, that is,
every quotient of two polynomials of the form (4.10), is continuous on $\R^k$
wherever the denominator is different from zero.\kn{An assembly line, and the
economical way to prove continuity ever after: two atoms (constants and
coordinates, the latter with the cleanest possible modulus,
$\delta = \varepsilon$) and three combining rules (sum, product, quotient
from 4.9; composition from 4.7). Almost no continuity proof after this page
touches $\varepsilon$ again --- functions are built out of continuous parts.
Chapter 8 extends the reach to power series, where the sums become
infinite and uniformity must be paid for.}

From the triangle inequality one sees easily that
\begin{equation}
  \bigl||\mathbf{x}|-|\mathbf{y}|\bigr| \le |\mathbf{x}-\mathbf{y}|
  \qquad (\mathbf{x},\mathbf{y} \in \R^k). \tag{4.11}
\end{equation}
Hence the mapping $\mathbf{x} \to |\mathbf{x}|$ is a continuous real function
on $\R^k$.

If now $\mathbf{f}$ is a continuous mapping from a metric space $X$ into
$\R^k$, and if $\phi$ is defined on $X$ by setting
$\phi(p) = |\mathbf{f}(p)|$, it follows, by Theorem 4.7, that $\phi$ is a
continuous real function on $X$.

\ritem{4.12}{Remark}
We defined the notion of continuity for functions defined on a subset $E$ of
a metric space $X$. However, the complement of $E$ in $X$ plays no role
whatever in this definition (note that the situation was somewhat different
for limits of functions). Accordingly, we lose nothing of interest by
discarding the complement of the domain of $f$. This means that we may just
as well talk only about continuous mappings of one metric space into another,
rather than of mappings of subsets. This simplifies statements and proofs of
some theorems. We have already made use of this principle in Theorems 4.8 to
4.10, and will continue to do so in the following section on
compactness.\kd{This is 2.16's observation --- every subset is a metric space
--- earning its keep, and it is subtler than it looks. Continuity of
$f|_E$ is measured \emph{inside} $E$: the function $f(x) = 1$ for rational
$x$, $0$ otherwise, is wildly discontinuous on $\R$, yet its restriction to
$\Q$ is constant, hence continuous. ``Continuous on $E$'' in Rudin always
means continuous as a function on the metric space $E$ --- the same
relativity as ``open in $E$'' (2.29--2.30), managed the same way: name the
space.}
